\documentclass[11pt,a4paper]{article}

% --- Packges ---
\usepackage[utf8]{inputenc}
\usepackage[margin=1in]{geometry}
\usepackage{graphicx}
\usepackage{amsmath,amssymb}
\usepackage{booktabs}
\usepackage{caption}
\usepackage{fancyhdr}
\usepackage{hyperref}
\usepackage{microtype}
\usepackage{multicol}
\usepackage{tikz}
\usetikzlibrary{shapes.geometric, arrows.meta}

% --- Custom Colors & Settings ---
\hypersetup{
    colorlinks=true,
    linkcolor=blue,
    filecolor=magenta,      
    urlcolor=cyan,
    pdftitle={TRANSIT - Project Report},
    pdfpagemode=FullScreen,
}

% --- Header/Footer Setup ---
\pagestyle{fancy}
\fancyhf{}
\fancyhead[L]{BAH 2026 - Problem Statement 7}
\fancyhead[R]{Team Bharat Ka Khazana}
\fancyfoot[C]{\thepage}
\renewcommand{\headrulewidth}{0.4pt}
\renewcommand{\footrulewidth}{0.4pt}

% ==========================================
%            DOCUMENT START
% ==========================================
\begin{document}

% ==========================================
%             PAGE 1: TITLE PAGE
% ==========================================
\begin{titlepage}
    \centering
    \vspace*{1.5cm}
    
    % Event Badge
    {\large\scshape BAH 2026 --- Problem Statement 7}\\[0.5cm]
    \rule{\linewidth}{1.5pt}\\[0.6cm]
    
    % Main Title
    {\Huge \bfseries TRANSIT: AI-Powered Exoplanet Detector}\\[0.3cm]
    {\large \bfseries AI-enabled Detection of Exoplanets from Noisy Astronomical Light Curves}\\[0.6cm]
    \rule{\linewidth}{1.5pt}\\[1.5cm]
    
    % Team Info
    {\large \bfseries Developed by Team: Bharat Ka Khazana}\\[1.0cm]
    
    % Members List
    \begin{minipage}{0.8\textwidth}
        \centering
        \begin{tabular}{r l}
            \textbf{Kritika Benjwal} & Lead AI Engineer \& Pipeline Architect \\
            \textbf{Sarthak Gupta} & Astrophysics \& Signal Processing Lead \\
            \textbf{Chaitanya Yadav} & Data Engineer \& Deployment Specialist \\
        \end{tabular}
    \end{minipage}\\[2.5cm]
    
    % Problem Statement Section
    \begin{minipage}{0.9\textwidth}
        \centering
        \textbf{\large Problem Statement Summary}\\[0.3cm]
        \textit{Astronomical surveys such as the Transiting Exoplanet Survey Satellite (TESS) observe millions of stellar targets, producing high-cadence light curves. Identifying planet candidates requires detecting minuscule, periodic brightness dips. However, stellar blending in crowded fields, instrumental anomalies, stellar spots, and astrophysical mimics (e.g. eclipsing binaries) heavily contaminate the signal. The goal is to build an end-to-end, robust automated pipeline that ingests, cleans, searches, classifies, and refines exoplanetary transits with high confidence.}
    \end{minipage}
    
    \vfill
    
    % Date
    {\large July 2026}
\end{titlepage}

\newpage
\clearpage

% ==========================================
%       PAGE 2: EXECUTIVE SUMMARY & METHODOLOGY
% ==========================================
\section{Executive Summary \& Objectives}
The discovery of exoplanets via the transit method relies on detecting periodic dips in stellar brightness when a planet passes in front of its host star. However, target stars located in crowded fields are subject to significant contamination. The **TRANSIT** pipeline is an automated, physics-informed deep learning system designed to extract and validate these signals. The core objectives are:
\begin{enumerate}
    \item \textbf{High-fidelity Ingestion:} Automated query and download of TESS 2-minute SPOC cadence light curves.
    \item \textbf{Astronomic Preprocessing:} Systematic removal of cosmic ray spikes and long-term stellar trend removal.
    \item \textbf{Physics-Based Detection:} Periodogram analysis utilizing Transit Least Squares (TLS) to model limb-darkened physical transit shapes.
    \item \textbf{Deep Learning Classification:} Categorizing signals into \textbf{Transit}, \textbf{Eclipsing Binary (EB)}, and \textbf{Other} (noise/systematics).
    \item \textbf{Parameter Refinement:} Real-time regression estimation of exoplanet parameters (period, depth, duration).
\end{enumerate}

\section{Pipeline Methodology}
The TRANSIT data pipeline operates as a sequential hybrid model combining classical signal analysis with neural network inference.

\begin{figure}[h]
    \centering
    \begin{tikzpicture}[nodeDistance=1.5cm, auto]
        % Styles
        \tikzStyle{block} = [rect, draw, fill=blue!10, text width=6cm, text centered, rounded corners, minimum height=1cm]
        \tikzStyle{arrow} = [thick,->,>=stealth]
        
        % Nodes
        \node [block] (tic) {TESS Catalog TIC ID};
        \node [block, below of=tic, node distance=1.3cm] (raw) {MAST Raw Data Download (Lightkurve)};
        \node [block, below of=raw, node distance=1.3cm] (clean) {MAD Sigma Clipping \& Wotan Detrending};
        \node [block, below of=clean, node distance=1.3cm] (tls) {Transit Least Squares (TLS) Period Search};
        \node [block, below of=tls, node distance=1.3cm] (fold) {Phase-Folding (200-Bin Profile \& 8-Bin Scalars)};
        \node [block, below of=fold, node distance=1.3cm] (model) {CNN-Transformer PyTorch Network};
        \node [block, below of=model, node distance=1.4cm] (out) {Classification (Transit/EB/Other) \\ \& Parameter Regression};
        
        % Arrows
        \draw [arrow] (tic) -- (raw);
        \draw [arrow] (raw) -- (clean);
        \draw [arrow] (clean) -- (tls);
        \draw [arrow] (tls) -- (fold);
        \draw [arrow] (fold) -- (model);
        \draw [arrow] (model) -- (out);
    \end{tikzpicture}
    \caption{Architecture overview of the TRANSIT exoplanet detection pipeline.}
    \label{fig:pipeline_flow}
\end{figure}

\subsection{Data Retrieval \& Preprocessing}
Using the \texttt{lightkurve} library, the pipeline queries the Mikulski Archive for Space Telescopes (MAST) and downloads target light curves. Instrumental outliers (e.g. focus changes, thruster firings) are removed using a three-step iterative \textbf{Median Absolute Deviation (MAD)} sigma-clipping algorithm set to limits of $4.0\sigma$ lower and $5.0\sigma$ upper. Stellar rotation and variability are flattened using a robust **biweight filter** via the \texttt{wotan} package with a window length of 0.5 days.

\subsection{Periodic Matched Filtering (Transit Least Squares)}
Unlike the traditional Box Least Squares (BLS) algorithm, which represents transits as flat-bottomed square boxes, the **Transit Least Squares (TLS)** algorithm is optimized for physical transit shapes. It incorporates realistic stellar limb-darkening parameters to model the ingress and egress transitions. TLS outputs:
\begin{itemize}
    \item \textbf{Signal Detection Efficiency (SDE):} Measures the peak prominence in the periodogram.
    \item \textbf{Signal-to-Noise Ratio (SNR):} Signifies transit depth divided by out-of-transit scatter.
    \item \textbf{Initial Physical Parameters:} Estimates period ($P$), transit depth ($\delta$), and duration ($d$).
\end{itemize}

\newpage
\clearpage

% ==========================================
%       PAGE 3: DEEP LEARNING, RESULTS & CHECKLIST
% ==========================================
\subsection{Hybrid CNN-Transformer Model Architecture}
The core classifier is a multi-modal PyTorch neural network that takes two distinct inputs:
\begin{enumerate}
    \item \textbf{Spatial Profile Branch (CNN):} A 1D convolutional neural network processes a 200-bin phase-folded transit profile to identify the local structure of the dip, extracting ingress/egress slopes and symmetry.
    \item \textbf{Self-Attention Sequence Block (Transformer):} Applies multi-head self-attention on positional encodings of the CNN feature maps. This captures global dependencies across the folded profile and highlights features critical for classification.
    \item \textbf{Metadata Parameter Branch (MLP):} A dense multi-layer perceptron processes 8 scalar features generated by TLS (SDE, SNR, period, transit count, transit depth, duration, secondary depth, and secondary flags).
\end{enumerate}
The spatial and scalar feature maps are fused and passed through a multi-task head: a Softmax classifier outputs probabilities for \textit{Transit}, \textit{Eclipsing Binary (EB)}, and \textit{Other}, while a regression head predicts corrected parameters.

\begin{table}[h]
    \centering
    \caption{Summary of core libraries and packages utilized in the project.}
    \begin{tabular}{lll}
        \toprule
        \textbf{Component} & \textbf{Library/Framework} & \textbf{Function} \\
        \midrule
        Data Access & \texttt{lightkurve} & Querying MAST archive and TESS metadata \\
        Preprocessing & \texttt{wotan} & Stellar variability detrending via biweight filter \\
        Matched Search & \texttt{transitleastsquares} & Physics-based limb-darkened periodogram search \\
        Deep Learning & \texttt{PyTorch} & Hybrid CNN-Transformer classification and regression \\
        Orbital Overlay & \texttt{batman-package} & Analytical light curve generation for validation \\
        User Interface & \texttt{Streamlit} & Dashboard compilation and UI layout rendering \\
        \bottomrule
    \end{tabular}
    \label{tab:libraries}
\end{table}

\section{Uncertainty, Significance \& Assumptions}
\subsection{Key Assumptions}
\begin{itemize}
    \item \textbf{Circular Orbits:} The pipeline assumes a circular orbit ($e=0$) during periodic matched searching.
    \item \textbf{Limb-Darkening Coefficients:} A quadratic limb-darkening model is used during visual validation with standard parameters ($u=[0.3, 0.2]$) for G-type stars.
\end{itemize}

\subsection{Significance Quantification}
To mitigate false positives, significance is verified through multiple levels:
\begin{enumerate}
    \item \textbf{Physical Indicators:} Signals with $SDE > 7.0$ and $SNR > 5.0$ are processed as strong planet candidates.
    \item \textbf{Statistical Confidence:} The network's Softmax output provides a posterior probability ($0.0 - 1.0$), which measures the classification certainty.
    \item \textbf{Cross-Validation:} The model was trained using 5-fold cross-validation, achieving an average validation accuracy of $\sim 95\%$.
\end{enumerate}

\section{Pipeline Validation Figures}
The figures below demonstrate the pipeline validation outputs generated during runtime:

\begin{figure}[h]
    \centering
    % In LaTeX, these lines will render the images when compiled locally with the files in your workspace directory.
    \begin{minipage}{0.48\textwidth}
        \centering
        \includegraphics[width=\textwidth]{results_eda/transit_data/phase_fold_profiles.png}
        \caption{Phase-folded light curve profiles.}
        \label{fig:phase_fold}
    \end{minipage}
    \hfill
    \begin{minipage}{0.48\textwidth}
        \centering
        \includegraphics[width=\textwidth]{models/transit_model/attention_heatmaps.png}
        \caption{Transformer self-attention heatmap.}
        \label{fig:attention_map}
    \end{minipage}
\end{figure}

\section{Project Checklist}
\begin{itemize}
    \item[\checkmark] \textbf{Light Curve Acquisition:} Automated TESS raw data download from MAST repository.
    \item[\checkmark] \textbf{Periodic Dip Search:} TLS periodogram analysis yielding SDE and SNR estimates.
    \item[\checkmark] \textbf{Classification Framework:} Multi-class classification (Transit / EB / Other) via CNN-Transformer.
    \item[\checkmark] \textbf{Parameter Estimation:} Physical parameter retrieval and neural network regression.
    \item[\checkmark] \textbf{Validation \& Dashboard:} Full visualization of curves, attention heatmaps, and analytical batman overlays.
\end{itemize}

\end{document}
